\documentclass[11pt]{article}
\usepackage{nopageno}
\usepackage[margin=1in]{geometry}
\usepackage{titlesec}
\titleformat{\section}[runin]{\normalfont\bfseries}{\thesection}{1em}{}[:]
\titlespacing*{\section}{0pt}{\baselineskip}{1ex}
\begin{document}
\pagestyle{empty}

\begin{center}
\textbf{Postdoctoral Mentoring Plan}
\end{center}

\vspace{-0.5cm}

\noindent This project supports two postdoctoral scholars: one jointly mentored by Napolitano and Busse (Penn State, spatial data science and hydrodynamic modeling) and one in Chi's lab (Indiana University, mixed-methods social science and network analysis). Because the project couples engineering and social science, each postdoc will develop skills in their home discipline while building working literacy in the other.
The goal of this plan is to prepare each postdoctoral scholar for an
independent research career by providing structured technical training,
professional development, and cross-disciplinary experience.
\vspace{-0.5cm}
\section*{Individual Development Plans and Performance Reviews} Within the first month of appointment, each postdoc will complete an
Individual Development Plan (IDP) with their supervising PI identifying
short- and long-term research goals, career goals, and skills to develop.
IDPs will be reviewed and updated annually. Each PI will conduct a formal
annual performance review assessing research progress, professional
development, and career goal progress. IDP completion and review will be
certified in NSF annual reports.
\vspace{-0.5cm}
\section*{Research Mentorship and Integration} Each PI will meet individually with their postdoc weekly to discuss research
progress, resolve technical challenges, and assess alignment with project
milestones. Monthly cross-PI video meetings bring all three postdocs together
with the full team, giving each postdoc regular exposure to the other
disciplinary perspectives central to the project. Each postdoc will lead
one research objective subtask, building ownership of a defined project
component while contributing to the broader integrated effort.
\vspace{-0.5cm}
\section*{Communication and Dissemination} Postdocs will present research at lab meetings, department seminars, and at
least one national conference annually, with travel supported by the project.
Target venues include the Engineering Mechanics Institute, the Association
for Preservation Technology annual conference, and the Natural Hazards
Research and Applications Workshop, reflecting the project's interdisciplinary
scope. Each PI will work directly with their postdoc to draft, submit, and
revise at least one peer-reviewed article per year, with the postdoc serving
as first author on work they lead. Postdocs will also contribute to
practitioner-facing toolkit materials and webinars, developing skills in
communicating technical findings to non-specialist audiences.
\vspace{-0.5cm}
\section*{Grant Writing and Research} Each postdoc will participate in the preparation of any grant renewals,
supplements, or new proposals submitted during the project period. PIs will
support postdocs in developing at least one independent funding application
during their appointment, such as an NSF postdoctoral fellowship or
relevant foundation award, providing feedback on drafts and connecting
postdocs with institutional grant development resources at their home
institution.
\vspace{-0.5cm}
\section*{Teaching, Mentoring, and Professional Skills} Each postdoc will have at least one structured opportunity per year to
mentor an undergraduate researcher, developing supervision and teaching
skills. Postdocs are encouraged to guest lecture in PI courses and to
attend workshops on pedagogy offered through their institution. All postdocs
will complete required training in responsible conduct of research and data
management practices consistent with NHERI DesignSafe open-data requirements.
Postdocs will be encouraged to attend professional development workshops on
topics including career planning, interview preparation, and academic job
market navigation.
\vspace{-0.5cm}
\section*{Cross-Disciplinary and Collaborative Training} Given the project's integration of engineering and social science, each
postdoc will attend at least one workshop or short course outside their home
discipline during the project period. PIs will facilitate introductions to
collaborators at the APT Disaster Response Initiative, USACE, and the Natural
Hazards Center to build postdocs' professional networks beyond their immediate
research community.
\vspace{-0.5cm}
\section*{Evaluation} The success of this mentoring plan will be assessed through annual IDP
reviews, postdoc self-evaluations, and tracking of career outcomes following
the appointment. PIs will solicit confidential feedback from postdocs on
mentoring quality and adjust plans accordingly.

\end{document}